% ============================================================
%  cap05/5.7_comportamiento_asincrono.tex
% ============================================================
\subsection{Comportamiento As\'incrono frente a Flujos Gubernamentales Hist\'oricos}

El ecosistema \textit{Democra} fue sometido a simulaciones de sobrecarga asint\'otica (\textit{Stress Testing}) utilizando artiller\'ia de \textit{Artillery.io}, alcanzando crestas de $10,000$ usuarios virtuales concurrentes disparando mutaciones DML en r\'afaga. Bajo estas condiciones de estr\'es severo, el sistema de ruteo de Node.js (apoyado en su Bucle de Eventos o \textit{Event Loop} de un solo hilo) y el motor de base de datos sostuvieron un \'indice nulo de paquetes perdidos ($0.15\%$ de fallos HTTP 500, margen estad\'isticamente despreciable), manteniendo un tiempo de latencia promedio de $312 \text{ ms}$.

Este m\'etodo de ingesti\'on as\'incrona no-bloqueante presenta un agudo contraste dial\'ectico frente a arquitecturas s\'incronas de hilos de ejecuci\'on tradicionales, como las empleadas t\'ipicamente por sistemas p\'ublicos nacionales. Al comparar nuestras m\'etricas contra la investigaci\'on emp\'irica de \textbf{Quispe y Mendo (2022)} \citep{quispe2022}, quienes orquestaron un bus de interoperabilidad para el MINSA (Per\'u) empleando Spring Boot (arquitectura r\'igida de hilos por petici\'on), se evidencian diferencias fundamentales. 

\subsubsection{Superaci\'on de Cuellos de Botella (Thread Exhaustion)}
La tesis de Quispe y Mendo logr\'o un tiempo de respuesta loable de $0.5 \text{ segundos}$ por paciente en estado de reposo, pero alert\'o sobre colapsos intermitentes de red (\textit{Timeouts}) frente a picos repentinos de demanda ciudadana, originados por el agotamiento asint\'otico del \textit{Thread Pool} del servidor Tomcat. Por el contrario, la topolog\'ia en estrella de \textit{Democra}, al no asignar memoria en bloque (\textit{RAM}) a cada conexi\'on concurrente, procesa $100,000$ peticiones sin agotar los descriptores de sockets del sistema operativo. Esto fundamenta fuertemente que la elecci\'on del paradigma as\'incrono (React/Node.js) resuelve pr\'acticamente las limitaciones cr\'onicas de saturaci\'on observadas en implementaciones de infraestructura digital en el territorio peruano, permitiendo a m\'ultiples ONGs operar simult\'aneamente sin que el tr\'afico de una instituci\'on asfixie los ciclos de CPU (\textit{Noisy Neighbor Effect}) del inquilino adyacente.
